% Algebra 2 Lesson 5-1 fast-track packet: nth Roots, Radicals, and Rational Exponents
% (enVision A2 2024 TE, Example 2 and Example 5 only; fast-lesson draft 2026-10-30, PLC has not met on this lesson)
% Bridge lesson (Precalculus and SAT).
% Quiz map (TE "5-1 Lesson Quiz"): item 1 real fourth roots (Example 1, cut; the +/- move is Example 5B);
% items 2-3 rational exponent <-> radical, evaluate (Example 2); item 4 simplify an nth root with
% variables (Example 4, cut); item 5 solve -4x^3 = 32 (Example 5).
% Shared preamble for the "fast" (hacked-down) Algebra 2 lesson packets, Oct 2026.
% Each packet: % Shared preamble for the "fast" (hacked-down) Algebra 2 lesson packets, Oct 2026.
% Each packet: % Shared preamble for the "fast" (hacked-down) Algebra 2 lesson packets, Oct 2026.
% Each packet: \input{fast-packet-template} then \begin{document} ... \end{document}.
\documentclass[12pt]{article}
\usepackage[margin=0.7in]{geometry}
\usepackage{amsmath,amssymb}
\usepackage{booktabs,tabularx,array}
\usepackage{enumitem}
\usepackage[T1]{fontenc}
\usepackage[utf8]{inputenc}
\usepackage{lmodern}
\usepackage{tikz}
\usepackage{pgfplots}
\pgfplotsset{compat=1.18}
\usepackage{xcolor}
\setlist{nosep,leftmargin=1.4em}
\setlength{\parindent}{0pt}
\setlength{\parskip}{6pt}
\renewcommand{\arraystretch}{1.35}
\newcolumntype{Y}{>{\raggedright\arraybackslash}X}
\pagestyle{empty}

% blanks and workspace
\newcommand{\blank}[1][2.6cm]{\rule[-2pt]{#1}{0.5pt}}
\newcommand{\lines}[1][1]{\par\vspace{2pt}\foreach \i in {1,...,#1}{\rule{0pt}{22pt}\rule{\linewidth}{0.4pt}\par}}
\newcommand{\wspart}[2]{\vspace{4pt}{\large\bfseries #1: #2}\par\vspace{-6pt}\rule{\textwidth}{0.4pt}\par}

% header: \packethead{Lesson 1-3}{Piecewise-Defined Functions}
\newcommand{\packethead}[2]{%
{\LARGE\bfseries #2}\hfill Name \blank[5cm]\\[2pt]
Algebra 2 \quad #1 (fast track) \hfill Date \blank[2.5cm] \quad Period \blank[1.2cm]\par\vspace{2pt}}

% boxed goals: \goals{learning goal}{language goal}
\newcommand{\goals}[2]{%
\fbox{\parbox{\dimexpr\textwidth-2\fboxsep-2\fboxrule}{%
\textbf{Learning goal:} #1\\[2pt]
\textbf{Language goal:} #2}}\par}

% quiz target box: \quiztarget{text}
\newcommand{\quiztarget}[1]{%
\vspace{4pt}\colorbox{gray!12}{\parbox{\dimexpr\textwidth-2\fboxsep}{\textbf{On the quiz:} #1}}\par}

% IXL table row helper: code & skill & which quiz question it feeds
\newcommand{\ixlrow}[3]{\texttt{#1} & #2 & #3 \\}

% one grid style for every graph; override xmin/xmax/... per plot
\pgfplotsset{sheetgrid/.style={
  axis lines=middle, grid=both, grid style={gray!45},
  xmin=-10.5, xmax=10.5, ymin=-10.5, ymax=10.5,
  xtick={-10,-8,...,10}, ytick={-10,-8,...,10},
  x=0.28cm, y=0.28cm,
  ticklabel style={font=\scriptsize, fill=white, inner sep=1pt},
  xlabel={$x$}, ylabel={$y$},
  axis line style={thick}, clip=true}}
% unit grid for small windows
\pgfplotsset{unitgrid/.style={sheetgrid,
  xmin=-5.5, xmax=5.5, ymin=-5.5, ymax=5.5,
  xtick={-5,...,5}, ytick={-5,...,5}, x=0.5cm, y=0.5cm}}
% open and closed endpoints
\tikzset{closed/.style={circle, fill, inner sep=1.6pt},
         open/.style={circle, draw, fill=white, inner sep=1.4pt}}
 then \begin{document} ... \end{document}.
\documentclass[12pt]{article}
\usepackage[margin=0.7in]{geometry}
\usepackage{amsmath,amssymb}
\usepackage{booktabs,tabularx,array}
\usepackage{enumitem}
\usepackage[T1]{fontenc}
\usepackage[utf8]{inputenc}
\usepackage{lmodern}
\usepackage{tikz}
\usepackage{pgfplots}
\pgfplotsset{compat=1.18}
\usepackage{xcolor}
\setlist{nosep,leftmargin=1.4em}
\setlength{\parindent}{0pt}
\setlength{\parskip}{6pt}
\renewcommand{\arraystretch}{1.35}
\newcolumntype{Y}{>{\raggedright\arraybackslash}X}
\pagestyle{empty}

% blanks and workspace
\newcommand{\blank}[1][2.6cm]{\rule[-2pt]{#1}{0.5pt}}
\newcommand{\lines}[1][1]{\par\vspace{2pt}\foreach \i in {1,...,#1}{\rule{0pt}{22pt}\rule{\linewidth}{0.4pt}\par}}
\newcommand{\wspart}[2]{\vspace{4pt}{\large\bfseries #1: #2}\par\vspace{-6pt}\rule{\textwidth}{0.4pt}\par}

% header: \packethead{Lesson 1-3}{Piecewise-Defined Functions}
\newcommand{\packethead}[2]{%
{\LARGE\bfseries #2}\hfill Name \blank[5cm]\\[2pt]
Algebra 2 \quad #1 (fast track) \hfill Date \blank[2.5cm] \quad Period \blank[1.2cm]\par\vspace{2pt}}

% boxed goals: \goals{learning goal}{language goal}
\newcommand{\goals}[2]{%
\fbox{\parbox{\dimexpr\textwidth-2\fboxsep-2\fboxrule}{%
\textbf{Learning goal:} #1\\[2pt]
\textbf{Language goal:} #2}}\par}

% quiz target box: \quiztarget{text}
\newcommand{\quiztarget}[1]{%
\vspace{4pt}\colorbox{gray!12}{\parbox{\dimexpr\textwidth-2\fboxsep}{\textbf{On the quiz:} #1}}\par}

% IXL table row helper: code & skill & which quiz question it feeds
\newcommand{\ixlrow}[3]{\texttt{#1} & #2 & #3 \\}

% one grid style for every graph; override xmin/xmax/... per plot
\pgfplotsset{sheetgrid/.style={
  axis lines=middle, grid=both, grid style={gray!45},
  xmin=-10.5, xmax=10.5, ymin=-10.5, ymax=10.5,
  xtick={-10,-8,...,10}, ytick={-10,-8,...,10},
  x=0.28cm, y=0.28cm,
  ticklabel style={font=\scriptsize, fill=white, inner sep=1pt},
  xlabel={$x$}, ylabel={$y$},
  axis line style={thick}, clip=true}}
% unit grid for small windows
\pgfplotsset{unitgrid/.style={sheetgrid,
  xmin=-5.5, xmax=5.5, ymin=-5.5, ymax=5.5,
  xtick={-5,...,5}, ytick={-5,...,5}, x=0.5cm, y=0.5cm}}
% open and closed endpoints
\tikzset{closed/.style={circle, fill, inner sep=1.6pt},
         open/.style={circle, draw, fill=white, inner sep=1.4pt}}
 then \begin{document} ... \end{document}.
\documentclass[12pt]{article}
\usepackage[margin=0.7in]{geometry}
\usepackage{amsmath,amssymb}
\usepackage{booktabs,tabularx,array}
\usepackage{enumitem}
\usepackage[T1]{fontenc}
\usepackage[utf8]{inputenc}
\usepackage{lmodern}
\usepackage{tikz}
\usepackage{pgfplots}
\pgfplotsset{compat=1.18}
\usepackage{xcolor}
\setlist{nosep,leftmargin=1.4em}
\setlength{\parindent}{0pt}
\setlength{\parskip}{6pt}
\renewcommand{\arraystretch}{1.35}
\newcolumntype{Y}{>{\raggedright\arraybackslash}X}
\pagestyle{empty}

% blanks and workspace
\newcommand{\blank}[1][2.6cm]{\rule[-2pt]{#1}{0.5pt}}
\newcommand{\lines}[1][1]{\par\vspace{2pt}\foreach \i in {1,...,#1}{\rule{0pt}{22pt}\rule{\linewidth}{0.4pt}\par}}
\newcommand{\wspart}[2]{\vspace{4pt}{\large\bfseries #1: #2}\par\vspace{-6pt}\rule{\textwidth}{0.4pt}\par}

% header: \packethead{Lesson 1-3}{Piecewise-Defined Functions}
\newcommand{\packethead}[2]{%
{\LARGE\bfseries #2}\hfill Name \blank[5cm]\\[2pt]
Algebra 2 \quad #1 (fast track) \hfill Date \blank[2.5cm] \quad Period \blank[1.2cm]\par\vspace{2pt}}

% boxed goals: \goals{learning goal}{language goal}
\newcommand{\goals}[2]{%
\fbox{\parbox{\dimexpr\textwidth-2\fboxsep-2\fboxrule}{%
\textbf{Learning goal:} #1\\[2pt]
\textbf{Language goal:} #2}}\par}

% quiz target box: \quiztarget{text}
\newcommand{\quiztarget}[1]{%
\vspace{4pt}\colorbox{gray!12}{\parbox{\dimexpr\textwidth-2\fboxsep}{\textbf{On the quiz:} #1}}\par}

% IXL table row helper: code & skill & which quiz question it feeds
\newcommand{\ixlrow}[3]{\texttt{#1} & #2 & #3 \\}

% one grid style for every graph; override xmin/xmax/... per plot
\pgfplotsset{sheetgrid/.style={
  axis lines=middle, grid=both, grid style={gray!45},
  xmin=-10.5, xmax=10.5, ymin=-10.5, ymax=10.5,
  xtick={-10,-8,...,10}, ytick={-10,-8,...,10},
  x=0.28cm, y=0.28cm,
  ticklabel style={font=\scriptsize, fill=white, inner sep=1pt},
  xlabel={$x$}, ylabel={$y$},
  axis line style={thick}, clip=true}}
% unit grid for small windows
\pgfplotsset{unitgrid/.style={sheetgrid,
  xmin=-5.5, xmax=5.5, ymin=-5.5, ymax=5.5,
  xtick={-5,...,5}, ytick={-5,...,5}, x=0.5cm, y=0.5cm}}
% open and closed endpoints
\tikzset{closed/.style={circle, fill, inner sep=1.6pt},
         open/.style={circle, draw, fill=white, inner sep=1.4pt}}

% Local header override: the TE title is too long to share a line with Name, so the title takes its own
% line and Name/Date/Period sit on the second line (template macro unchanged). "(fast track)" is dropped so
% the Name blank keeps the template's 5 cm; the Period blank matches the template.
\renewcommand{\packethead}[2]{%
{\LARGE\bfseries #2}\par\vspace{4pt}
Algebra 2 \quad #1 \hfill Name \blank[5cm] \quad Date \blank[2cm] \quad Period \blank[1.2cm]\par\vspace{2pt}}
\begin{document}

\packethead{Lesson 5-1}{nth Roots, Radicals, and Rational Exponents}
\goals{I can relate rational exponents to roots, evaluate expressions with rational exponents, and solve equations by taking the nth root of both sides.}{I can explain what a rational exponent means, and how I took a root to solve an equation, using the words root, power, index and radicand.}
\vspace{-4pt}{\footnotesize Spanish: \emph{índice} (index), \emph{radicando} (radicand), \emph{exponente} (exponent), \emph{símbolo radical} (radical symbol).}

\vspace{-4pt}\quiztarget{(1) the real fourth roots of a positive number: two, $\pm$ (full Savvas quiz only; Example 5B). (2) pick every form equal to a rational exponent and (3) say what a rational exponent means, then evaluate it: Example 2, Try It 2, Practice. (4) simplify an $n$th root with variables (full Savvas quiz only; Example 4, IXL LQX). (5) solve $ax^3=c$: Example 5, Try It 5. Questions 3 and 5 count 10 points each; 2 is the +5 bonus.}

\wspart{Example 2}{Understand Rational Exponents}
\begin{minipage}[t]{0.5\textwidth}
\vspace{0pt}\setlength{\abovedisplayskip}{3pt}\setlength{\belowdisplayskip}{3pt}
\textbf{A.} What is the meaning of the exponent in the expression $16^{\frac14}$?

Extend the properties of integer exponents.
\begin{align*}
16^{\frac14}&=\left(2^4\right)^{\frac14} && 16=2^4\\
&=2^{4\cdot\frac14}=2^1 && \text{Power of a Power}\\
&=2
\end{align*}
Since $2^4=16$, $16^{\frac14}$ is $\sqrt[4]{16}$, the fourth \textbf{root} of 16, which is 2.

In general ($n$ a positive integer), $x^{\frac1n}=\sqrt[n]{x}$, the principal $n$th root: one value, real when $x\ge0$ or $n$ is odd.
\end{minipage}\hfill
\begin{minipage}[t]{0.46\textwidth}
\vspace{0pt}\setlength{\abovedisplayskip}{3pt}\setlength{\belowdisplayskip}{3pt}
\textbf{B.} Write $27^{\frac23}$ as a radical.
\begin{align*}
27^{\frac23}&=\left(27^{\frac13}\right)^2 && \text{Power of a Power}\\
&=\left(\sqrt[3]{27}\right)^2\\
&=3^2=9
\end{align*}
27 is a perfect cube ($3^3$), so take the cube root first, then square: $\left(\sqrt[3]{27}\right)^2=9$.

Same value the other way (power first, then root):\\ $\sqrt[3]{27^2}=\sqrt[3]{729}=9$.
\end{minipage}

\vspace{2pt}
\fbox{\parbox{\dimexpr\textwidth-2\fboxsep-2\fboxrule}{%
\textbf{The rule.} \ When $\sqrt[n]{c}$ is real, $c^{\frac mn}=\left(\sqrt[n]{c}\right)^m=\sqrt[n]{c^m}$.\\
Denominator $n$ $=$ the \textbf{index} (which \textbf{root}); \ numerator $m$ $=$ the \textbf{power}; \ $c$ $=$ the \textbf{radicand}.}}

\textbf{Say it:} In $27^{\frac23}$, the denominator 3 is the \blank[2cm] of the radical, so I take the cube \blank[1.6cm]. The numerator 2 is the \blank[1.8cm], so I square. The \blank[2.2cm] is 27.

\wspart{Try It 2}{Explain what each fractional exponent means, then evaluate}
\begin{minipage}[t]{0.48\textwidth}
\vspace{0pt}
\textbf{a.} $25^{\frac12}$

$25^{\frac12}$ is the \blank[2.4cm] root of 25.
\lines[2]
Value: \blank[2cm]
\end{minipage}\hfill
\begin{minipage}[t]{0.48\textwidth}
\vspace{0pt}
\textbf{b.} $32^{\frac25}$

$32^{\frac25}$ is the \blank[2.4cm] root of 32, \blank[2.2cm].
\lines[2]
Value: \blank[2cm]
\end{minipage}

\newpage
% R5 Practice: rehearses Lesson Quiz item 2 (select all expressions equivalent to x^(3/4)), which no Try It covers.
% Structural parallel only: same select-all structure (three equivalent forms, three wrong forms), different exponent 2/5.
\wspart{Practice}{Which forms are equal?}
Select \textbf{all} the expressions that are equivalent to $y^{\frac25}$. Use the rule: denominator $=$ index, numerator $=$ power.

\vspace{2pt}
\begin{tabular}{@{}l@{\hspace{2.2em}}l@{\hspace{2.2em}}l@{}}
$\square$ \ \textbf{A.} $\sqrt[5]{y^2}$ & $\square$ \ \textbf{B.} $\left(y^{\frac15}\right)^2$ & $\square$ \ \textbf{C.} $\sqrt{y^5}$ \\[8pt]
$\square$ \ \textbf{D.} $y^{\frac{10}{4}}$ & $\square$ \ \textbf{E.} $\left(y^2\right)^{\frac15}$ & $\square$ \ \textbf{F.} $\left(y^{\frac12}\right)^{\frac15}$
\end{tabular}

\vspace{2pt}
For one choice you did \emph{not} select, write it as a single power of $y$ and say what went wrong.
\lines[1]

\wspart{Example 5}{Use nth Roots to Solve Equations}
% TE Example 5 callout prints (x^n)^(1/n) = x^1 for every n; that holds for odd n only (even n gives |x|), so the idea is split by index.
\textbf{The idea:} get $x^n$ alone, then take the $n$th \textbf{root} of both sides: raise each side to the $\frac1n$ \textbf{power}. Odd~$n$: $\left(x^n\right)^{\frac1n}=x$, one solution. Even~$n$: $\left(x^n\right)^{\frac1n}=|x|$, so $x^n=d$ with $d>0$ gives $x=\pm d^{\frac1n}$ (no real solution if $d<0$).

\begin{minipage}[t]{0.5\textwidth}
\vspace{0pt}
\textbf{A.} Solve $2x^5=64$.
\begin{align*}
2x^5&=64\\
x^5&=32 && \text{divide by 2}\\
\left(x^5\right)^{\frac15}&=32^{\frac15} && \text{raise to the } \tfrac15\\
x&=2 && 2^5=32
\end{align*}
The graphs of $y=2x^5$ and $y=64$ meet once, so there is only one real solution: $x=2$.
\end{minipage}\hfill
\begin{minipage}[t]{0.46\textwidth}
\vspace{0pt}
\begin{tikzpicture}
\begin{axis}[sheetgrid, xmin=-5.4, xmax=5.4, ymin=-12, ymax=82,
  xtick={-4,-2,2,4}, ytick={20,40,60}, minor xtick={-5,...,5}, minor ytick={-10,0,...,80},
  grid=both, minor grid style={gray!45}, x=0.58cm, y=0.035cm, yticklabel style={anchor=north east}]
\addplot[very thick, red, domain=-1.3:2.08, samples=80, <->] {2*x^5};
\addplot[very thick, blue, domain=-5.2:5.2, samples=2, <->] {64};
\node[font=\scriptsize, fill=white, inner sep=1pt, anchor=north west] at (axis cs:0.1,-1) {$O$};
\node[closed] at (axis cs:2,64) {};
\node[font=\scriptsize, red, fill=white, inner sep=1pt, anchor=west] at (axis cs:2.4,40) {$y=2x^5$};
\node[font=\scriptsize, blue, fill=white, inner sep=1pt, anchor=south] at (axis cs:-3.5,66) {$y=64$};
\end{axis}
\end{tikzpicture}
\end{minipage}

\begin{minipage}[t]{0.5\textwidth}
\vspace{0pt}
\textbf{B.} Solve $3x^4=1875$.
\begin{align*}
3x^4&=1875\\
x^4&=625 && \text{divide by 3}\\
% TE Example 5B prints (x^4)^(1/4) = +/-(625)^(1/4); the left side is |x|, so the packet writes x = +/-(625)^(1/4).
x&=\pm\left(625\right)^{\frac14} && \text{even index, } 625>0\text{: } \pm\\
x&=\pm5 && 5^4=625
\end{align*}
$(-5)^4=625$ too: a negative number to an even power is positive. The graphs meet twice, so the solutions are $x=5$ and $x=-5$.
\end{minipage}\hfill
\begin{minipage}[t]{0.46\textwidth}
\vspace{0pt}
\begin{tikzpicture}
\begin{axis}[sheetgrid, xmin=-10.6, xmax=10.6, ymin=-250, ymax=2650,
  xtick={-8,-4,4,8}, ytick={400,800,1200,1600,2000,2400}, yticklabels={400,800,1200,,2000,2400},
  minor xtick={-10,-8,...,10}, minor ytick={-200,0,...,2600},
  grid=both, minor grid style={gray!45}, x=0.29cm, y=0.00116cm, yticklabel style={anchor=south east}]
\addplot[very thick, red, domain=-5.4:5.4, samples=80, <->] {3*x^4};
\addplot[very thick, blue, domain=-10.3:10.3, samples=2, <->] {1875};
\node[font=\scriptsize, fill=white, inner sep=1pt, anchor=north west] at (axis cs:0.2,-20) {$O$};
\node[closed] at (axis cs:-5,1875) {};
\node[closed] at (axis cs:5,1875) {};
\node[font=\scriptsize, red, fill=white, inner sep=1pt, anchor=west] at (axis cs:5.6,900) {$y=3x^4$};
\node[font=\scriptsize, blue, fill=white, inner sep=1pt, anchor=south] at (axis cs:-7.5,1950) {$y=1875$};
\end{axis}
\end{tikzpicture}
\end{minipage}

\vspace{2pt}
{\footnotesize
\begin{tabularx}{\textwidth}{|Y|Y|Y|}
\hline
\textbf{$x^n=c$, index $n$ odd} & \textbf{index $n$ even, $c>0$} & \textbf{index $n$ even, $c<0$} \\ \hline
one real root, same sign as $c$: \ $x^3=-8$, \ $x=-2$ & two real roots, $\pm$: \ $x^4=625$, \ $x=\pm5$ & no real root: \ $x^4=-256$ \\ \hline
\end{tabularx}}

\textbf{Say it:} To solve $3x^4=1875$, I divided both sides by \blank[1cm], then took the fourth \blank[1.6cm] of both sides (raised them to the $\frac14$ \blank[1.6cm]). The \blank[1.6cm] 4 is even and the \blank[2cm] 625 is positive, so there are \blank[1cm] real solutions.

\newpage
\wspart{Try It 5}{Solve each equation}
\begin{minipage}[t]{0.48\textwidth}
\vspace{0pt}
\textbf{a.} Solve the equation $5x^3=320$.
\lines[6]
$x=$ \blank[2.5cm]\\[4pt] How many real solutions? \blank[1cm]
\end{minipage}\hfill
\begin{minipage}[t]{0.48\textwidth}
\vspace{0pt}
\textbf{b.} Solve the equation $2p^4=162$.
\lines[6]
$p=$ \blank[2.5cm]\\[4pt] How many real solutions? \blank[1cm]
\end{minipage}

\vspace{6pt}
Why does part b have more solutions than part a? Use the words \emph{index} and \emph{even}.\lines[2]

\wspart{IXL}{Practice for the quiz}
Go to IXL, Algebra 2. Type the code in the search box. Do the starred skills first.

\vspace{2pt}
\begin{tabularx}{\textwidth}{l Y Y}
\toprule
\textbf{Code} & \textbf{Skill} & \textbf{Feeds} \\
\midrule
\ixlrow{9TM $\bigstar$}{Convert between rational exponents and radicals}{Quiz 2, 3 (Example 2)}
\ixlrow{KJX $\bigstar$}{Evaluate rational exponents}{Quiz 3 (Example 2, Try It 2)}
\ixlrow{U42 $\bigstar$}{Nth roots}{Quiz 5 (Example 5); Quiz 1, $\pm$ (full Savvas quiz only)}
\ixlrow{EUH}{Roots of integers}{Warm-up, optional: the perfect powers 32, 243, 625}
\ixlrow{LQX}{Simplify radical expressions with variables I}{Quiz 4 (full Savvas quiz only), optional}
\bottomrule
\end{tabularx}

\textbf{SmartScore 70 or above on each starred ($\bigstar$) skill replaces your quiz score.}

\newpage
% R3: the lesson's only DOK 3 (Practice & Problem Solving 50, Performance Task). The TE item analysis files it
% under Example 6 (cut), which is Example 5's x^n = c move in a context, so it is named as the self-paced extension.
\wspart{Extension (self-paced)}{Practice 50 (DOK 3, Performance Task)}
\fbox{\parbox{\dimexpr\textwidth-2\fboxsep-2\fboxrule}{%
A milk processing company uses cylindrical-shaped containers. The height of the container is equal to the diameter of the base. The volume of a container is $169.65\text{ ft}^3$.

\vspace{2pt}
\textbf{Part A} How much material is needed to make the lateral surface area of the shipping container?

\textbf{Part B} The cargo hold of a ship is 20 ft high. What is the largest number of these shipping containers that could be stacked on top of each other inside the cargo hold?

\vspace{4pt}
{\small Hint: write $V$ using $r$ only, then solve $r^3=c$ as in Example 5.}
\lines[9]
Part A: \blank[3cm] ft$^2$ \qquad Part B: \blank[2cm] containers\par\vspace{2pt}}}

\textbf{Before the quiz, finish these sentences.}

In $c^{\frac mn}$, the denominator $n$ is the \blank[2cm], the numerator $m$ is the \blank[2cm], and $c$ is the \blank[2.4cm].

To solve $ax^n=c$, I divide by $a$, then take the $n$th \blank[2cm] of both sides. If the index is even and $\frac ca>0$, I write \blank[1cm] (two real solutions); if the index is even and $\frac ca<0$, there is \blank[1.6cm] real solution.

\end{document}
